\begin{textsample}{POS Dim 3 – 2018\_c5 – Score 1.00 – 2018\_c5/enade\_2018\_secr\_exec\_32b5ea331f.txt}  \label{ex:f3_pos_182}
CAPÍTULO 5 DISTRIBUIÇÃO DOS CONCEITOS 5.1 PANORAMA NACIONAL DA DISTRIBUIÇÃO DOS CONCEITOS O Conceito Enade é um indicador de qualidade que avalia os cursos por intermédio dos desempenhos dos estudantes a partir dos resultados do Enade.
Os valores possíveis vão de 1 ( pior situação ) a 5 ( melhor situação ).
Os cursos com apenas um ou sem nenhum concluinte participante não obtêm o Conceito Enade, ficando Sem Conceito ( SC ).
Na Tabela 5.1, são apresentadas a quantidade e a distribuição de cursos de Secretariado Executivo participantes do Enade/2018, por faixa de conceito e Grande Região.
Enfatiza -se, mais uma vez, que a diferença entre os cursos tabulados neste capítulo e no capítulo 2 corresponde aos cursos Sem Conceito, ou seja àqueles, em princípio, sem alunos concluintes que houvessem participado da prova ou àqueles que tivessem tão somente um aluno concluinte19.
Observando -se os dados da Tabela 5.1, nota -se que, dos 42 cursos participantes, 15 ( 35,7 \% ) classificaram -se com conceito 3, o valor modal.
Esse foi também o conceito modal em duas regiões : Sudeste ( 35,7 \% ) e Sul ( 52,9 \% ).
O conceito 2 foi o segundo mais frequente em nível nacional ( 23,8 \%, correspondendo a dez cursos ), e os conceitos 1 e 4, os terceiros ( 11,9 \%, correspondendo a cinco cursos, cada ).
Houve, ainda, quatro cursos ( 9,5 \% ) que receberam conceito 5.
Dos 42 cursos de Secretariado Executivo, três cursos ficaram Sem Conceito ( SC ).
A região Norte participou com dois cursos ou 4,8 do total nacional.
Desses, um curso recebeu o conceito 1, o que equivale a 50,0 \% do total regional, e um outro recebeu o conceito 19 Esses cursos com somente um concluinte não foram considerados no capítulo 2. 2.
Nenhum curso ficou Sem Conceito ( SC ), e nenhum curso recebeu os demais conceitos restantes.
A região Nordeste participou com cinco cursos ou 11,9 \% do total nacional.
Nessa região, dois cursos ( 40,0 \% em termos regionais ) obtiveram conceito 2, o conceito modal para a região.
Os conceitos 3, 4 e 5 foram atribuídos, respectivamente, a um curso cada ( 20,0 \%, cada ).
Nenhum curso ficou Sem Conceito ( SC ) ou recebeu conceito 1.
Dos 14 cursos participantes da região Sudeste ( 33,3 \% do total nacional ), o conceito 3 foi o conceito modal, como já anteriormente assinalado, tendo sido obtido por cinco cursos ( 35,7 \% ).
O conceito 2 foi atribuído a três cursos ( 21,4 \% ) e o conceito 4, a outros três.
E somente um curso ( 7,1 \% ) obteve conceito 5.
Nessa região, dois cursos ficaram Sem Conceito ( SC ) e nenhum recebeu o conceito 1.
Os 17 cursos da região Sul corresponderam a 40,5 \% do total nacional.
Como já indicado, houve predominância do conceito 3, atribuído a 52,9 \% dos cursos da região ( conceito modal ), ou seja, atribuído a nove dos 17 cursos participantes na região Sul.
Seis cursos ( 11,8 \% ) receberam conceitos 1, 2 e 5, respectivamente, dois cursos, receberam o conceito 1, outros dois, o conceito 2, e outros 2, o conceito 5 ( 11,8 \%, cada ).
E apenas um curso ( 5,9 \% ) obteve conceito 4.
Um curso ficou Sem Conceito ( SC ) na região Sul.
Já dos quatro cursos participantes na região Centro-Oeste ( 9,5 \% do total nacional ), dois ( 50,0 \% em termos regionais, cada ) receberam conceito 1 e outros dois, o conceito 2.
Nenhum curso ficou Sem Conceito ( SC ) ou recebeu os demais conceitos restantes nessa região.
No Gráfico 5.1, é apresentada a distribuição cumulativa do Conceito Enade segundo Grande Região.
Em um gráfico de distribuição cumulativa, em cada valor de abscissa ( eixo x ), o valor apresentado no eixo das ordenadas ( eixo y ) é o valor acumulado da variável.
Por exemplo, nas regiões Norte ( linha azul ) e Centro-Oeste ( linha \textbf{laranja} ), linhas sobrepostas, todos os cursos apresentam conceito 2 ou 1.
A região Nordeste ( linha verde ), por outro lado, apresenta 60 \% dos cursos com conceito 3 ou abaixo.
Nas regiões Sudeste ( linha vermelha ) e Sul ( linha rosa ), 91,7 \% e 87,5 \% dos cursos, respectivamente, têm conceito 4 ou abaixo disso.
Em linhas gerais, em um gráfico de distribuição cumulativa, poligonais mais à esquerda representam Grandes Regiões com pior distribuição, e poligonais mais à direita, Grandes Regiões com melhores conceitos.
As regiões Norte ( linha azul ) e Centro-Oeste ( linha \textbf{laranja} ) apresentam os cursos com conceitos mais baixos, ao passo que as regiões Nordeste ( linha verde ), Sudeste ( linha vermelha ) e Sul ( linha rosa ) apresentam as distribuições com valores maiores. 5.2 CONCEITOS POR CATEGORIA ADMINISTRATIVA, POR MODALIDADE DE ENSINO E GRANDE REGIÃO Na Tabela 5.2, é apresentada a distribuição dos cursos de Secretariado Executivo participantes do Enade/2018, por Categoria Administrativa e Modalidade de ensino, de acordo com os conceitos por eles alcançados, segundo as Grandes Regiões brasileiras.
Dos 42 cursos participantes, 18 ( 42,9 \% ) eram ministrados em Instituições Públicas, e 24 ( 57,1 \% ), em Instituições Privadas.
Quanto à Modalidade de Ensino, apenas um ( 2,4 \% ) curso era a Distância, e 41 ( 97,6 \% ), Presenciais.
De acordo com as informações da Tabela 5.2, em termos nacionais, entre Instituições Públicas e Privadas, observa -se que, dos quatro cursos avaliados com conceito 5, todos eram oferecidos em IES Públicas.
Dos 18 cursos participantes de IES Públicas, o conceito 3 foi o valor modal, atribuído a seis cursos ( 33,3 \% ).
Dos demais cursos de Instituições Públicas participantes, quatro cursos foram avaliados com conceito 2, enquanto outros quatro obtiveram conceito 5, como já comentado e dois cursos ( 11,1 \% ) foram avaliados com conceito 1, enquanto outros dois receberam conceito 4.
Nenhum curso ficou Sem Conceito ( SC ) nesta categoria.
Na rede privada, também, o conceito modal foi 3, atribuído a nove cursos, o correspondente a 37,5 \% dos 24 cursos da categoria.
Dos demais cursos participantes, três ( 12,5 \% ) receberam conceito 1, enquanto outros três obtiveram conceito 4 e seis ( 25,0 \% ), conceito 2.
Nesta Categoria Administrativa, três cursos ficaram Sem Conceito ( SC ) e nenhum curso recebeu conceito 5.
Dos 41 cursos participantes da modalidade presencial, o conceito 3 foi o valor modal, atribuído a 14 cursos ( 34,1 \% ).
Dos demais cursos nessa Modalidade de Ensino, dez cursos ( 12,2 \% ) receberam conceito 1 ou 4 ( cinco cursos cada ), dez cursos ( 24,4 \% ) receberam conceito 2, e quatro cursos ( 9,8 \% ), conceito 5.
Nessa Modalidade de Ensino, três cursos ficaram Sem Conceito ( SC ).
Ao único curso a Distância informado atribuiu -se o conceito 3.
Nessa Modalidade de ensino, tampouco, nenhum curso ficou Sem Conceito ( SC ).
Na análise por região, observa -se que, na região Norte, as Instituições Privadas não tiveram participação com nenhum curso.
As Instituições Públicas participaram com dois cursos na região Norte ( totalidade regional ), a um dos quais foi atribuído conceito 1, enquanto um outro curso recebeu conceito 2.
Nenhum curso ficou Sem Conceito ( SC ) ou com conceito 3, 4 ou 5 nessa combinação de categoria e região.
Todos os cursos oferecidos na região Norte eram presenciais.
Na região Nordeste, a rede privada não obteve nenhum curso participante.
As Instituições Públicas da região Nordeste participaram com cinco cursos ( totalidade da região ).
Desses, dois obtiveram conceito 2, o valor modal, seguidos por outros três cursos, que receberam conceito 3, 4 ou 5 ( um curso cada ).
Nenhum curso ficou Sem Conceito ( SC ) ou recebeu conceito 1 nessa região.
Todos os cursos oferecidos na região Nordeste eram presenciais.
Na região Sudeste, a proporção de cursos da rede privada, 92,9 \%, foi mais elevada do que a registrada nas demais regiões brasileiras, correspondendo a 13 dos 14 cursos participantes.
Nessa categoria e região, o conceito modal foi 3, atribuído a cinco cursos.
Três cursos ( 23,1 \% ) obtiveram conceito 2, enquanto outros três receberam conceito 4.
Dois cursos ficaram Sem Conceito ( SC ) e nenhum curso recebeu os conceitos 1 ou 5.
Somente um curso foi oferecido por Instituições Públicas na região Sudeste, e seu conceito foi 5.
Nessa região, nenhum curso de IES Públicas ficou Sem Conceito ( SC ) ou recebeu os demais conceitos restantes.
Todos os cursos oferecidos na região Sudeste eram presenciais.
As Instituições Privadas concentraram oito dos 17 cursos participantes da região Sul, 47,1 \% do total regional.
Desses, quatro cursos obtiveram conceito 3, o conceito modal.
Nesta combinação de Categoria Administrativa e Grande Região, dois cursos receberam conceito 2 e um curso, conceito 1.
Um curso ficou Sem Conceito ( SC ) e nenhum curso obteve conceitos 4 ou 5.
As Instituições Públicas na região Sul participaram com nove cursos ( 52,9 \% dos cursos da região ).
Desses, cinco receberam conceito 3, o conceito modal.
Dos demais cursos, um recebeu conceito 1, e um outro, conceito 4, e dois cursos obtiveram conceito 5.
Nenhum curso da região Sul oferecido por IES Pública ficou Sem Conceito ( SC ) ou recebeu o conceito 2.
Na região Sul, a maioria dos cursos foi oferecido na modalidade Presencial : 16 cursos ( 94,1 \% ).
Para esses, a moda foi o conceito 3, com oito cursos.
Quanto aos demais cursos presenciais da região, dois foram alocados ao conceito 1, outros dois receberam o conceito 2, enquanto ainda outros dois obtiveram conceito 5, e somente um curso recebeu o conceito 4.
Um curso ficou Sem Conceito ( SC ) nessa região.
Na região Sul, o único curso oferecido, na modalidade Educação a Distância, obteve o conceito 3.
Na região Centro-Oeste, três dos quatro cursos participantes eram de Instituições Privadas ( 75,0 \% em termos regionais ).
Desses, dois receberam conceito 1, o conceito modal, e um curso, o conceito 2.
Nenhum curso ficou Sem Conceito ( SC ), e nenhum curso recebeu os demais conceitos restantes.
Apenas um curso foi oferecido por Instituições Públicas na região Centro-Oeste ( 25,0 \% do total regional ), e ele foi alocado ao conceito 2.
Nenhum curso ficou Sem Conceito ( SC ) ou recebeu os demais conceitos nessa combinação de categoria e região.
Todos os cursos oferecidos na região Centro-Oeste eram presenciais.
No Gráfico 5.2, apresenta -se a distribuição do Conceito Enade, segundo a Categoria Administrativa da IES.
Os cursos em IES Públicas ( linha azul ) apresentam uma poligonal mais à direita do que os em IES Privadas ( linha verde ) e, consequentemente, uma distribuição de Conceitos Enade com valores maiores. 5.3 CONCEITOS POR ORGANIZAÇÃO ACADÊMICA E GRANDE REGIÃO Na Tabela 5.3, encontra -se a distribuição dos conceitos atribuídos aos cursos participantes do Enade/2018, na Área de Secretariado Executivo, por Organização Acadêmica, segundo as Grandes Regiões brasileiras.
Dos 42 cursos de Secretariado Executivo participantes, 27 eram oferecidos em Universidades, sete, em Centros Universitários, sete, em Faculdades e um, em CEFET/IFET.
Esta distribuição corresponde a, respectivamente, 64,3 \%, 16,7 \%, 16,7 \% e 2,4 \% dos cursos.
De acordo com os dados apresentados, de todos os quatro cursos avaliados com conceito 5, todos eram vinculados a Universidades.
Os cursos vinculados a esse tipo de Organização Acadêmica tiveram o conceito 3 como conceito modal, atribuído a nove cursos ( 33,3 \% ).
Os demais cursos vinculados a Universidades receberam os conceitos 1 ( três cursos ), 2 ( seis cursos ), 4 ( quatro cursos ) e conceito 5 ( quatro cursos, como já mencionado ).
Um curso ficou Sem Conceito ( SC ) nessa categoria.
Dos cursos em Centros Universitários, o conceito modal também foi 3, atribuído a quatro cursos ( 57,1 \% ).
Dois cursos vinculados a esse tipo de Organização Acadêmica receberam os conceitos 2 ou 4 ( um curso, cada ).
Nessa categoria, um ficou Sem Conceito ( SC ) ou recebeu os conceitos 1 ou 5.
Dos seis cursos mantidos por Faculdades, dois ( 28,6 \% ) foram atribuídos ao conceito 1, outros dois receberam conceito 2, enquanto outros dois receberam conceito 3.
Dos cursos oferecidos em Faculdades, um ficou Sem Conceito ( SC ) e nenhum curso recebeu os conceitos 4 ou 5.
O único curso mantido por CEFET/IFET, obteve conceito 2.
Considerando -se, separadamente, as regiões brasileiras, verifica -se que, na região Norte, as Universidades participaram com a totalidade dos cursos avaliados ( dois cursos ).
A um curso foi atribuído conceito 1, enquanto um outro obteve conceito 2.
Nenhum curso ficou Sem Conceito ( SC ) ou recebeu os demais conceitos restantes nessa combinação de categoria e região.
Os Centros Universitários, as Faculdades e os CEFET/IFET da região Norte não foram representados por nenhum curso.
Na região Nordeste, as Universidades participaram, também, com a totalidade de cursos da Área de Secretariado Executivo oferecidos ( cinco cursos ).
Nenhum desses cinco cursos ficou Sem Conceito ( SC ).
O conceito modal 2 foi atribuído a dois cursos.
Os três cursos restantes obtiveram conceitos 3, 4 e 5, respectivamente.
Nenhum curso oferecido por Universidades no Nordeste recebeu conceito 1.
Assim como na região Norte, os Centros Universitários, as Faculdades e os CEFET/IFET da região Nordeste não foram representados por nenhum curso Na região Sudeste, as Universidades concentraram sete dos 14 cursos de Secretariado Executivo da região.
Dos cursos oferecidos em Universidades, nessa região, os conceitos modais foram 2 e 4, atribuídos a dois cursos cada.
Dos demais cursos, um recebeu o conceito 3, enquanto um outro obteve conceito 5.
Um curso ficou Sem Conceito ( SC ) e nenhum curso recebeu o conceito 1.
Os Centros Universitários participaram com três cursos na região Sudeste, dos quais dois obtiveram o conceito modal 3.
O conceito 4 foi atribuído a um curso.
Nenhum curso ficou Sem Conceito ( SC ) ou recebeu os demais conceitos restantes nessa combinação de organização acadêmica e região.
Já as Faculdades, foram representadas por quatro cursos na região Sudeste.
O conceito modal, também 3, foi atribuído a dois cursos.
Já o conceito 2 foi atribuído a um curso.
Um curso ficou Sem Conceito ( SC ) e nenhum curso recebeu os demais conceitos restantes nessa categoria.
Os CEFET/IFET não tiveram participação com nenhum curso nessa combinação de organização acadêmica e região.
Dos 17 cursos da região Sul, 12 eram mantidos por Universidades, para os quais o conceito modal foi 3, atribuído a sete cursos.
Os demais cursos receberam os conceitos 1 e 5 ( dois cursos cada ).
Um curso recebeu conceito 4.
Nesse tipo de organização, nenhum curso da região Sul ficou Sem Conceito ( SC ) ou recebeu o conceito 2.
Dos quatro cursos participantes de Centros Universitários da região Sul, o conceito modal, também 3, foi atribuído a dois cursos, e um curso obteve o conceito 2.
Um curso ficou sem Conceito ( SC ), e nenhum curso recebeu os demais conceitos restantes nessa combinação de organização acadêmica e região.
Somente um curso foi vinculado a Faculdades na região Sul, e esse recebeu conceito 2.
Nenhum curso ficou Sem Conceito ( SC ) ou recebeu os demais conceitos restantes.
Já os CEFET/IFET não participaram com nenhum curso na região Sul.
Na região Centro-Oeste, um dos quatro cursos oferecidos era mantido por Universidades.
Para este tipo de Organização Acadêmica, esse curso ficou alocado ao conceito 2.
Nesse tipo de organização, nenhum curso ficou Sem Conceito ( SC ) ou recebeu os demais conceitos restantes.
Os Centros Universitários da região Centro-Oeste não contaram com nenhum curso.
Dos dois cursos oferecidos em Faculdades na região Centro-Oeste, receberam o conceito 1.
Nesta região, nenhum curso ficou Sem Conceito ( SC ), e nenhum obteve os demais conceitos restantes.
Já os CEFET/IFET participaram com apenas um curso na região Centro-Oeste, ao qual foi atribuído o conceito 2 nessa combinação de organização acadêmica e região.
No Gráfico 5.3, apresenta -se a distribuição do Conceito Enade segundo a Organização Acadêmica da IES.
O único curso em CEFET/IFET ( linha \textbf{laranja} ) recebeu nota 2 e é uma poligonal em escada, e não entrará na análise que se segue.
Os cursos em Universidades ( linha azul ) e Centros Universitários ( linha verde ) apresentam poligonais mais à direita do que as demais, denotando uma melhor distribuição de Conceitos Enade.
No outro extremo, encontram -se os cursos mantidos por Faculdades ( linha vermelha ) apresentando a respectiva poligonal mais à esquerda e, consequentemente, distribuição de Conceitos Enade com menores valores.

% matched lemmas: laranja
\end{textsample}
